\section{Resultados y controles}
\label{sec:resultados}

Todos los números de esta sección provienen de los archivos de resultados del cálculo (\texttt{resultados/PT\_k.npz}, \texttt{resultados/checks.json}) o de las funciones de fondo del código, y las figuras se generan a partir de esos archivos con scripts versionados (Apéndice~\ref{ap:codigo}). Los parámetros son los de la tabla~\ref{tab:decisiones}. Salvo que se diga lo contrario, las cantidades dimensionales están en unidades $\MP=1$, $m=1$.

\subsection{El fondo: dónde y por qué difieren RG y UG}
\label{sec:fondores}

Con las mismas condiciones iniciales ($\phi_i=25.386\,\MP$, $H_i=10.369\,m$, $\epsilon_1^{\rm ini}=0.0031$ en los dos modelos), la evolución posterior es muy distinta (figura~\ref{fig:fondo}). En RG la inflación dura $161.83$ e-folds y termina cuando $\phi$ llega a $\phi\simeq1.01\,\MP$ (en slow-roll estricto sería $\sqrt2\,\MP\simeq1.41\,\MP$; la diferencia refleja que aquí $\epsilon_1=1$ se evalúa con la dinámica completa). En UG, con $Q_i=0.1\,V(\phi_i)=32.22$ y $\gamma=0.1$, la duración fijada es de $100.00$ e-folds. El parámetro de Hubble de UG decrece más rápido que el de RG, y $\epsilon_1$ de UG es mayor en todo el intervalo.

\begin{figure}[t]
\centering
\includegraphics[width=\textwidth]{fig4_fondo.pdf}
\caption{Fondo cosmológico en RG (azul) y UG con difusión (naranja), con las mismas condiciones iniciales. (a) $H(N)$; (b) $\epsilon_1(N)$; (c) $\phi(N)$; (d) presupuesto de energía en UG: $3H^2=\rho_\phi+Q(N)-Q_i$. Generada por \texttt{make\_report\_figures.py}.}
\label{fig:fondo}
\end{figure}

Hay dos hechos de esta figura que hay que entender, porque condicionan el resto.

El primero es el salto inicial de $\epsilon_1$ de UG, que en el primer e-fold pasa de $0.0031$ a $0.0104$ (máximo en $N\simeq1.04$) y luego baja a un mínimo de $\sim0.006$. Se debe a que la velocidad inicial $\dot\phi_i$ es la de slow-roll de RG, que no está en el atractor de UG: el término $-\dot Q/\dot\phi$ de \eqref{eq:KGQ} es un empuje adicional que, en slow-roll, es del orden de $\gamma\,(Q/V)\,V^2/V'^2\times V'=\gamma\,(Q/V)\,(\phi^2/4)\,V'\simeq1.6\,V'$ al inicio (estimación de orden de magnitud). El campo se acelera hasta un nuevo régimen cuasi-estacionario en un e-fold. Es una consecuencia de imponer las mismas condiciones iniciales en ambos modelos, y no un error.

El segundo hecho es más importante, y no era evidente antes de graficar: \emph{el final de la inflación en UG no ocurre porque el campo llegue al fondo del potencial}. En $N=100$ el campo de UG está en $\phi=8.44\,\MP$, con $\rho_\phi=37.35$ y $\epsilon_V=2/\phi^2\simeq0.03$ (este epsilon del potencial desnudo no implica slow-roll del fondo UG, que tiene $\epsilon_1=1$; un universo de RG con ese mismo valor del campo tendría por delante unos $\phi^2/4-1/2\simeq17$ e-folds de inflación), pero $3H^2=5.13$. La razón está en el panel (d): como $Q(N)$ decae exponencialmente hasta $Q(100)=1.5\times10^{-3}$, la ecuación de Friedmann \eqref{eq:H2UG} con $\Lambda_0=-Q_i$ queda $3H^2=\rho_\phi-Q_i+Q\simeq\rho_\phi-32.2$: la constante de integración negativa le resta a $H^2$ una energía constante comparable a $\rho_\phi$ cuando este ha bajado a $\sim37$. Con $H^2$ tan pequeño, $\epsilon_1=\dot\phi^2/2H^2$ llega a $1$. Con más precisión: si $K=\dot\phi^2/2$ y $U=V+Q+\Lambda_0$ (en $\MP=1$), es $3H^2=K+U$ y $\epsilon_1=3K/(K+U)$, de modo que el criterio $\epsilon_1=1$ implementado equivale a $U=2K$, \emph{no} a $U=0$: al final, $3H^2=5.13$ da $K=1.71$ y $U=3.42>0$. La constante negativa $\Lambda_0=-Q_i$ resta a $U$ y adelanta ese instante respecto de RG, pero el final de la inflación no coincide con el cero del potencial efectivo. Es decir: en este modelo, la supresión de $H$ y el final temprano de la inflación en UG los produce en buena medida la elección $\Lambda_0=-Q_i$, que a su vez se sigue de pedir el mismo $H_i$ en ambos modelos. Esto se discute en la sección~\ref{sec:decisiones-dep}.

\subsection{Un modo a través del horizonte}

La figura~\ref{fig:modo} muestra la evolución de un modo individual con $k=2.12\times10^{18}\,m$, que sale del horizonte en $N=40.000$ en RG y en $N=40.149$ en UG. Antes del cruce, el modo oscila con amplitud constante (para $\tilde v$) y con oscilaciones que se van estirando en $N$ al acercarse al cruce, porque su frecuencia respecto de $N$ es $k/(aH)$ (panel b); es la oscilación de Bunch--Davies en un fondo en expansión. En el cruce, el $P_T$ instantáneo del modo, $2k^3|h|^2/\pi^2$, deja de decrecer y se congela (panel a). Los valores finales son $5.888\times10^{-10}$ en RG y $4.368\times10^{-10}$ en UG, con cociente $0.742$.

\begin{figure}[t]
\centering
\includegraphics[width=\textwidth]{fig5_modo.pdf}
\caption{Un modo con $k=2.12\times10^{18}\,m$ en RG (azul) y UG (naranja). (a) $P_T$ instantáneo del modo; se congela al cruzar el horizonte (líneas grises). (b) Parte real de $\tilde v_k$ antes del cruce. Generada por \texttt{make\_report\_figures.py}.}
\label{fig:modo}
\end{figure}

Este cociente puede anticiparse casi sin integrar. En $N=40$ los parámetros de Hubble son $H_{\rm RG}=8.994$ y $H_{\rm UG}=7.760$, de modo que $(H_{\rm UG}/H_{\rm RG})^2=0.744$, que reproduce el valor numérico $0.742$ con la diferencia que cabe esperar de las correcciones de slow-roll y del pequeño corrimiento del instante de cruce. Es exactamente lo que dice la sección~\ref{sec:espectro}: el espectro se fija en el cruce del horizonte y depende de $H$ en ese instante.

\subsection{El espectro $P_T(k)$}

La figura~\ref{fig:PTk} es el resultado principal del cálculo: el espectro tensorial en RG y UG a igual $k$ comóvil, para $40$ modos entre $k=1.05\times10^4\,m$ y $k=3.48\times10^{39}\,m$ (un rango de $81.8$ en $\ln k$, limitado por la duración de la inflación de UG), con el cociente debajo.

\begin{figure}[t]
\centering
\includegraphics[width=\textwidth]{fig1_PT_k_RG_vs_UG.pdf}
\caption{Espectro tensorial primordial en RG y UG con difusión $Q(N)$. Arriba, $P_T(k)$ con la diferencia sombreada; abajo, el cociente UG/RG. Los marcadores huecos son la fórmula slow-roll de RG evaluada con el fondo de UG (control C3.4). La banda gris marca los modos donde $\epsilon_1^{\rm UG}>0.02$ en el cruce. El eje superior reetiqueta los mismos $k$ en e-folds. Generada por \texttt{make\_figures.py} a partir de los resultados guardados.}
\label{fig:PTk}
\end{figure}

En RG, $P_T$ decrece suavemente de $7.49\times10^{-10}$ a $3.50\times10^{-10}$ a lo largo del rango (una pendiente ligeramente roja, $n_T\simeq-2\epsilon_1$). En UG, para los mismos modos, decrece de $6.85\times10^{-10}$ a $1.21\times10^{-10}$. El cociente UG/RG va de $0.914$ en el primer modo a $0.344$ en el último, y en los nueve modos que se usan en los controles vale $0.888$, $0.825$, $0.782$, $0.742$, $0.698$, $0.644$, $0.574$, $0.477$ y $0.327$. Para estos parámetros, entonces, UG suprime $P_T$ frente a RG a igual $k$ en una cantidad que va del 8.6\% al 65.6\% según el modo. La banda gris marca la zona en que $\epsilon_1$ de UG en el cruce supera $0.02$, es decir, los modos que salen cuando la inflación de UG está terminando y la aproximación de slow-roll deja de ser buena: la magnitud del cociente ahí refleja en buena parte cuánto cae $H$ hacia el final, como se explicó en la sección~\ref{sec:fondores}.

Un punto a no perder de vista: los espectros se comparan a igual $k$ comóvil, con el mismo $a_i$ y las mismas condiciones iniciales, y como las dos inflaciones tienen duraciones distintas ($161.8$ y $100$ e-folds), la correspondencia entre un $k$ dado y ``$N$ e-folds antes del final de la inflación'' \emph{no} es la misma en los dos modelos. Por eso las figuras no deben leerse como una comparación a igual escala observacional, que requeriría modelar el recalentamiento y fijar un $k_*$ observable en cada modelo. \emph{La sección~\ref{sec:potenciales} muestra que esto no es un detalle: a igual número de e-folds antes del final de cada inflación el signo de la diferencia se invierte en el cuadrático} ($P_T^{\rm UG}/P_T^{\rm RG}=1.5$ en lugar de $0.3$--$0.9$). El resultado de esta subsección, ``UG suprime $P_T$'', vale solo a igual $k$ comóvil.

\subsection{Controles}

Los resultados anteriores se someten a diez controles reproducibles (nueve de resultado y uno de control de mutaciones), cada uno con tolerancia fijada \emph{antes} de correrlo y que no se modificó al ver los resultados; el código es \texttt{checks.py} y la evidencia legible por máquina es \texttt{checks.json}. Se los agrupa según lo que pretenden mostrar.

\subsubsection{El límite estándar de RG}

El primer control usa un fondo artificial de de Sitter exacto ($H$ constante, $\epsilon_1=0$), donde el modo de Bunch--Davies es analítico y $P_T=\frac{2H^2}{\pi^2\MP^2}(1+(k/aH)^2)$. El código reproduce esa expresión con un error relativo de $8.7\times10^{-7}$ para $k=10^3$, $10^5$ y $10^8$ (constante en $k$, como debe ser por la invariancia de escala). La solución a tiempo finito se obtiene de \cite[ecs.~196 y 198]{baumann22}, con la normalización tensorial usada aquí. Es una prueba a la vez del integrador, de la condición de Bunch--Davies, de la normalización de $P_T$ (incluido el factor de la sección~\ref{sec:normalizacion}) y de la forma en $N$ de la ecuación de modos. El error coincide en magnitud con el que produce el arranque en $k/aH=100$ en el barrido de convergencia.

El segundo compara, en RG con $V=\tfrac12m^2\phi^2$, el $P_T$ numérico con el desarrollo de slow-roll de Martin, Ringeval y Vennin \cite{martin14} a orden cero, uno y dos en los parámetros de flujo. La tabla~\ref{tab:sr} muestra el residuo relativo de $P_T/(2H^2/\pi^2\,a_0)-1$ en cuatro modos. A orden cero, la desviación es $-2(C+1)\epsilon_1$ (que es lo que predice \eqref{eq:PTNLO}); cada orden reduce el residuo, y a segundo orden queda por debajo de $6\times10^{-7}$ para $N_{\rm exit}=20$ a $100$, es decir, del orden del error de la condición inicial. En $N_{\rm exit}=140$, donde $\epsilon_1=0.023$, el residuo es $-4\times10^{-5}$ (unas $3\epsilon_1^3$, el tamaño esperado del término de orden siguiente).

\begin{table}[t]
\centering\small
\renewcommand{\arraystretch}{1.2}
\begin{tabular}{@{}rcccc@{}}
\toprule
$N_{\rm exit}$ & $\epsilon_1$ & LO & NLO & NNLO\\
\midrule
20 & 0.0035 & $-1.92\times10^{-3}$ & $-1.15\times10^{-5}$ & $+4.9\times10^{-7}$\\
60 & 0.0049 & $-2.69\times10^{-3}$ & $-2.28\times10^{-5}$ & $+5.0\times10^{-7}$\\
100 & 0.0081 & $-4.45\times10^{-3}$ & $-6.38\times10^{-5}$ & $-2.9\times10^{-7}$\\
140 & 0.0230 & $-1.30\times10^{-2}$ & $-5.51\times10^{-4}$ & $-4.1\times10^{-5}$\\
\bottomrule
\end{tabular}
\caption{Residuo relativo de $P_T$ numérico (RG, $\phi^2$) frente al desarrollo de slow-roll a orden cero (LO), uno (NLO) y dos (NNLO), evaluado en el cruce del horizonte.}
\label{tab:sr}
\end{table}

El tercero es la comparación con las predicciones conocidas de $V\propto\phi^2$. Para $N_*=60$ e-folds antes del final, el cálculo da $\epsilon_1=0.00836$, una relación $r=P_T/P_\mathcal{R}^{\rm LO}=0.1332$ (usando $P_\mathcal{R}=H^2/8\pi^2\epsilon_1$ de \eqref{eq:PRr} con el $H$ y el $\epsilon_1$ del código), que es 99.5\% de $16\epsilon_1=0.1338$; y para $N_*=50$, $r=0.1597$. Están dentro del rango de ``aproximadamente $0.15$'' que Planck da para el potencial cuadrático \cite{planck20}. La pendiente tensorial, medida como diferencia finita en $\ln k$ entre modos, coincide con la de la fórmula de \cite{martin14} a $2\times10^{-6}$: $n_T=-0.016947$ frente a $-0.016945$ (y $-2\epsilon_1=-0.016727$ a orden más bajo). Finalmente, la amplitud escalar $A_s=P_\mathcal{R}^{\rm LO}m_{\rm phys}^2$ a $N_*=60$ es $2.167\times10^{-9}$, frente al valor de Planck $2.099\times10^{-9}$ ($\ln10^{10}A_s=3.044$): una diferencia de $+3.2$\%, que valida el valor $m=6\times10^{-6}\MP$ elegido. Este último control usa la fórmula estándar del espectro escalar de RG, no un cálculo propio.

\subsubsection{Convergencia numérica}

Se varía un parámetro numérico por vez y se compara con un valor de referencia más fino, en cuatro modos y en RG y UG (figura~\ref{fig:conv}). Con los valores por defecto el resultado es estable a $\sim10^{-6}$: las tolerancias de los integradores contribuyen menos de $10^{-9}$, la resolución de los splines del fondo no se distingue de $4\times10^{-13}$ a partir de $2\,001$ puntos, y las dos fuentes dominantes son la condición inicial (el arranque en $k/aH=100$ da $9.0\times10^{-7}$, y cae como $(aH/k)^3$: $3.6\times10^{-5}$ para $30$, $3.8\times10^{-8}$ para $300$) y el punto de medida ($10^{-6}$, el $(k/aH)^2$ esperado). Los cinco valores por defecto están dentro de su tolerancia a priori.

\begin{figure}[t]
\centering
\includegraphics[width=\textwidth]{fig2_convergencia.pdf}
\caption{Convergencia numérica de $P_T$: error máximo frente a un valor de referencia más fino, variando un parámetro por vez (cuatro modos, RG y UG). La línea es la tolerancia fijada a priori y el círculo marca el valor por defecto.}
\label{fig:conv}
\end{figure}

\subsubsection{La diferencia entre UG y RG}

\begin{figure}[t]
\centering
\includegraphics[width=\textwidth]{fig3_limite_estandar.pdf}
\caption{(A) RG con $\phi^2$: residuo de $P_T$ frente al desarrollo de slow-roll a tres órdenes. (B) El cociente UG/RG numérico frente al de la fórmula de slow-roll de RG evaluada con el fondo de UG.}
\label{fig:limite}
\end{figure}

Para que la diferencia de la figura~\ref{fig:PTk} sea un resultado, tiene que descartarse que sea un error numérico, y hay que entender de dónde viene. Se hicieron cinco comprobaciones (a las que se suma la de sensibilidad de la sección siguiente).

Primero, el ruido: recalculando el cociente con distintas tolerancias de los integradores, con otro punto de arranque y con otro punto de medida, cambia como máximo $8.5\times10^{-7}$ (dominado por el arranque), frente a una señal mínima de $0.112$: un cociente entre la diferencia de modelos y el error numérico estimado de $1.3\times10^{5}$ (es un cociente de tamaños numéricos, no una significancia estadística ni una detectabilidad experimental). Segundo, se compara el cociente numérico con el que da la fórmula de slow-roll de RG \eqref{eq:cociente} evaluada con el $H$, $\epsilon_1$ y $\epsilon_2$ \emph{propios de UG} en su cruce del horizonte (figura~\ref{fig:limite}B): coinciden a $\lesssim3\times10^{-5}$ en los ocho modos donde se aplica el criterio fijado a priori ($\epsilon_1$ y $|\epsilon_2|$ de ambos modelos menores que $0.05$), la mayoría a $\sim10^{-6}$, y en el noveno (con $\epsilon_2^{\rm UG}=0.10$, fuera del criterio) la diferencia es $2.7\times10^{-4}$. Tercero, un control de consistencia entre los dos códigos: con $\gamma=0$ ($Q$ constante y $\Lambda_0=-Q_i$), UG debe ser RG exacta (sección~\ref{sec:difusionQ}), y así ocurre a $4.3\times10^{-13}$ en $P_T$ y $8.5\times10^{-14}$ en la duración: prueba el camino de UG del código (\texttt{quantities\_ug}, la ecuación \eqref{eq:KGQ} con $\dot Q$ y \eqref{eq:Lambda0}) contra el de RG, que son funciones distintas. Cuarto, la ecuación \eqref{eq:HdotUG}, que no se usó para construir el fondo, se verifica sobre la solución integrada: $|\!-\!\dd\ln H/\dd N-\phi_N^2/2|$ es a lo sumo $7.9\times10^{-12}$ en RG y $3.0\times10^{-7}$ en UG (en $N=0$, donde $\epsilon_1$ de UG salta). Quinto, un integrador independiente de la ecuación \eqref{eq:modoX}, escrito en tiempo cósmico, para $h$ y sin usar $N$, $v$, $\epsilon_1$ ni splines, coincide con el de \texttt{modes.py} a $4.4\times10^{-10}$ en RG y UG; valida los pasos que llevan de \eqref{eq:modoX} a \eqref{eq:modoN}.

\subsubsection{Que los controles puedan detectar errores}

Un control que nunca falla no prueba nada, y por eso se agregan dos comprobaciones de sensibilidad. Un control negativo: si se calcula el modo de UG con un $\Lb$ equivocado (omitiendo $Q$, de modo que $X\neq0$ y aparece una masa efectiva $-2X$), $P_T$ cambia en un 60\% para $N_{\rm exit}=12$, 7.6\% para $40$, 0.7\% para $70$ y 0.3\% para $85$ (la sensibilidad decrece porque $Q\propto\ee^{-\gamma N}$). Con el $X$ que se calcula frente a forzar $X=0$, el cambio es menor que $5\times10^{-14}$: el término de masa calculado es despreciable. Y un control de mutaciones: se rompe el código a propósito de tres maneras (duplicar la normalización de $P_T$, cambiar el signo de $\epsilon_1$ en la ecuación de modos, y desplazar $\Lambda_0$ en 0.01\% de $Q_i$) y se verifica que los cuatro controles correspondientes fallan (el cambio de signo de $\epsilon_1$ debe ser detectado por dos de ellos).

\subsection{Balance de los controles}

Los diez controles pasan, y también la batería completa de pruebas automáticas (17 en el momento de esta sección, contando la sincronía del notebook con el código y las pruebas de las figuras; las secciones~\ref{sec:potenciales} y~\ref{sec:acoplamiento} agregan ocho de Starobinsky y cuatro simbólicas). Vale la pena decir qué muestran y qué no. Muestran que el código reproduce el límite estándar de RG con la precisión que se le pide, que su resultado es estable frente a los parámetros numéricos, y que la diferencia entre UG y RG es real y se explica por el fondo. \emph{No} muestran que la derivación analítica sea correcta: el control de atribución valida la numérica, y la ecuación de modos de UG que integra el código es, por construcción, la de la derivación con el término de masa que se demostró que se anula; el único control que podría contradecirla es el de sensibilidad (que dice que, si no se anulara, el código lo vería), pero éste no la comprueba de manera independiente. Tampoco hay una referencia externa contra la cual comparar $P_T$ de UG con difusión, porque no se conoce una en la literatura leída; UG se apoya en consistencia interna y con RG. Estos límites se retoman en la sección~\ref{sec:discusion}.
